% fisicai analysis note template — structure mirrors experiment publications
% (Introduction / Detector and data / Selection / Method / Systematics / Results / Summary).
% Copy into note/ with `python -m fisicai.writeup --init-note note/` and fill in.
\documentclass[11pt]{article}
\usepackage{hepnote}
% Auto-generated by `python -m fisicai.writeup` -- do not edit by hand.
\newcommand{\NEventsTotal}{83343}
\newcommand{\NEventsTwoMuons}{11983}
\newcommand{\NEventsOppositeCharge}{10420}
\newcommand{\NEventsMassWindow}{8589}
\newcommand{\BinWidthGev}{0.5}
\newcommand{\MassWindowLo}{60}
\newcommand{\MassWindowHi}{120}
\newcommand{\FitLo}{85}
\newcommand{\FitHi}{97}
\newcommand{\PtMin}{20}
\newcommand{\EtaMax}{2.4}
\newcommand{\MZPeak}{91.002}
\newcommand{\MZPeakStat}{0.041}
\newcommand{\MZPeakSystRange}{0.122}
\newcommand{\SigmaRes}{1.19}
\newcommand{\SigmaResStat}{0.07}
\newcommand{\NSigFit}{6446}
\newcommand{\NSigFitStat}{191}
\newcommand{\NBkgFit}{978}
\newcommand{\ChiTwo}{25.1}
\newcommand{\Ndf}{19}
\newcommand{\ChiTwoNdf}{1.32}
\newcommand{\GammaZPdg}{2.4955}
\newcommand{\MZPdg}{91.1880}
\newcommand{\MZPdgErr}{0.0020}
\newcommand{\PeakMinusPdgMev}{-186}
\newcommand{\PeakMinusPdgPermille}{2.0}


\title{\bfseries TITLE: physics object and measurement, dataset}
\author{fisicai analysis bundle}
\date{\today}

\begin{document}
\notenumber{AN-XXXX-XX}
\makenoteheader
\maketitle

\begin{abstract}
One paragraph: dataset and integrated luminosity or event count; selection in one
sentence; method; the result with \stat and \syst uncertainties quoted via macros;
one sentence of interpretation and its limits.
\end{abstract}

\section{Introduction}
Physics context with citations; what is determined here and why; the reproducibility
statement (every number is generated by \texttt{analysis.py} and imported as a macro).

\section{Detector and data}
One paragraph on the detector, cited. Dataset provenance (Open Data record, derived
format), run period, \sqs, and the branches/objects used.

\section{Event selection}
Selection criteria with motivation, then the cutflow:
\begin{table}[htbp]
  \centering
  \caption{Event selection and yields.}
  \begin{tabular}{lr}
    \toprule
    Selection & Events \\
    \midrule
    All events            & \NEventsTotal \\
    % ... one row per cut, macros only ...
    \bottomrule
  \end{tabular}
\end{table}

\section{Analysis method}
Estimator or fit model, fit range, what is fixed vs floated, and why.

\section{Systematic uncertainties}
What is evaluated (and how) and, explicitly, what is \emph{not} evaluated.
Use a table if more than two sources.

\section{Results}
The key figure (mplhep style, experiment-label conventions) and the result sentence
with \stat\syst uncertainties. Compare to external references in prose with citations.

\section{Summary}
Three sentences: what was done, the result, what it does and does not claim.

\bibliographystyle{unsrt}
\bibliography{references}
\end{document}
